\documentclass[11pt]{article}
\usepackage{amsmath,amssymb}
\title{Exact finite-boundary cancellation before nonlinear flux power counting}\author{}\date{}
\begin{document}\maketitle
\paragraph{The coefficient is a boundary functional.}
Let $\mathcal C_R(\eta)$ be the COMPLETE first nonlinear coefficient of
the original conditional compact partition function, as in the first original nonlinear flux coefficient calculation and the local gauge completion first nonlinear coefficient calculation.
The boundary record $\eta$ includes all fixed tangential boundary links.
The classical minimum, complete collars, full Haar factors and the
nonlinear gauge Jacobian are part of its definition. This functional
exists smoothly near the flat record for each fixed finite $R$.
No size-uniform expansion is assumed here.

For any gauge transformation $g$ on the boundary vertices, extend it to
the interior vertices in any way and change the interior integration
variables by the corresponding site gauge map. Product Haar is exactly
preserved and every original plaquette is conjugated. Thus, with the
actual fixed-only action still included,
\[
 Z_R(\eta^g)=Z_R(\eta),\qquad v_R(\eta^g)=v_R(\eta).
\]
Uniqueness of the fixed-$R$ asymptotic coefficients gives
\[
 \mathcal C_R(\eta^g)=\mathcal C_R(\eta).                           \tag{1}
\]
This uses the original finite integral, not an infinite-volume Ward
identity or an unproved diagram cancellation. Restoring the interior
tree gauge after the change is an exact interior gauge transformation;
it does not alter the boundary record or its Haar normalization.

\paragraph{A complete Hessian Ward identity.}
Use right boundary logarithms $b_e$ near identity. Global color
conjugation implies $\nabla\mathcal C_R(0)=0$: there is no invariant
linear functional on one adjoint $SU(2)$ index. Let $\Pi_R$ be the Hessian
of $\mathcal C_R$ at zero, including its actual minimum/frame response.
The infinitesimal boundary gauge generator has value $D_\partial\xi$
at zero. Differentiating (1) in an arbitrary boundary direction gives
\[
 \boxed{\quad \Pi_R D_\partial=0,\qquad
                D_\partial^*\Pi_R=0.\quad}                       \tag{2}
\]
The derivative of the nonlinear gauge generator multiplies
$\nabla\mathcal C_R(0)$ and hence vanishes. This is where the gauge
acceleration is retained and eliminated honestly, rather than ignored.
All quartic, cubic-pair, ghost, Haar and changing-covariance terms in
the connected first nonlinear flux response calculation/the local gauge completion first nonlinear coefficient calculation must be combined for (2). It need not hold for them
separately.

Restrict to color three and let $\mathcal B_1$ be the fixed boundary
links in direction one. The boundary gauge parameter $\xi_x=x_1i\sigma_3$
has constant boundary difference on $\mathcal B_1$ and zero on other
directions. Thus, for EVERY boundary row $e$,
\[
 \sum_{f\in\mathcal B_1}(\Pi_R)_{ef}=0.                            \tag{3}
\]
There are no missing opposite-face or corner rows: the parameter is
defined on the entire original boundary. For the direction-one block,
(3) fixes its diagonal local term by the complete off-diagonal sum.
It is the exact cancellation of the zero-momentum boundary mass term.
It does not assert that an isolated bulk diagram has zero mass term.

\paragraph{Actual flux response as a finite-boundary difference form.}
The physical constant-flux boundary source is
$a_e=x_0/R^2$ on $\mathcal B_1$, zero otherwise. Symmetry of $\Pi_R$
and (3) give the exact complete identity
\[
 \boxed{\quad
 \mathcal C_R''[a,a]
 =-\frac1{2R^4}\sum_{e,f\in\mathcal B_1}
    (\Pi_R)_{ef}\,(x_0(e)-x_0(f))^2.\quad}                       \tag{4}
\]
This removes the potential-coordinate amplitude before any absolute
Wick estimate. It retains a squared physical spatial difference and
all original finite-boundary terms. In particular shifting the Landau
origin $a_e\mapsto(x_0-c)/R^2$ leaves (4) unchanged for every $c$.
The complete boundary mass and the mixed constant/flux response vanish,
not merely their leading large-$R$ approximations.

An exact rational regression shifts the original $R=2$ flux origin
to $c=0,2,-3$, re-expands every original plaquette and all covariance,
quartic, cubic and Haar partners. In every case the complete
$\mathcal C_2''$ is exactly $-119/20736$.
The finite calculation records this
test. It checks the finite-boundary cancellation, not a growth law.

\paragraph{Scale accounting after the cancellation.}
The boundary link set has $O(R^3)$ points. For each boundary edge, the
number of boundary edges in a Euclidean shell of radius $r$ is at most
$Cr^2$ up to the usual unit-thickness convention; equivalently dyadic
balls contain at most $Cr^3$. This includes distinct faces and corners
with a fixed dimension-dependent multiplicity. Let
\[
 W_k=\sum_{\substack{e,f\in\mathcal B_1\\
                 2^{k-1}\le1+|e-f|<2^k}}|(\Pi_R)_{ef}|.
\]
The actual identity (4) gives
\[
 |\mathcal C_R''[a,a]|
       \le \frac{C}{R^4}\sum_{k\le C+\log_2R}2^{2k}W_k.           \tag{5}
\]
Thus a sufficient concrete next estimate is
\[
 W_k\le C R^3\,2^{-k}(1+k)^p.                                   \tag{6}
\]
It would yield $C(1+\log R)^p$, normalized by the bounded physical
energy $E_R$, since the dyadic sum is dominated by $R(1+\log R)^p$.
A pointwise boundary kernel bound
$|(\Pi_R)_{ef}|\le C(1+|e-f|)^{-4}
 (1+\log(2+|e-f|))^p$ would imply (6).
Neither bound is proved here. They concern the COMPLETE boundary kernel,
whose derivatives include harmonic/minimum response, not a bare bulk
vertex covariance. The boundary dimension is three; inserting a
four-dimensional shell count here would give the wrong target.

\paragraph{What has been gained and what remains.}
Equations (2)--(4) prove the dangerous complete zero-momentum
cancellation with the exact finite boundary and without taking
absolute values of the individual diagrams. They provide the
specific second spatial moment that a scale decomposition must
estimate. They do not bound that moment. Gauge symmetry permits
curvature and other physical boundary responses; it does not fix a
marginal coupling correction or imply (6). The complementary cancellation
of longitudinal propagators must
likewise precede, not replace, the remaining connected spatial estimate.
The original differentiated nonlinear remainder remains a separate
required bound. No uniform nonlinear S94 result is claimed.
\end{document}
